% Generated by manuscript/build_manuscript.py.
\setcounter{table}{0}
\setcounter{figure}{0}
\section{Experiment Details}
\label{app:experiment-details}


\subsection{Metric Details}


We use mean squared error (MSE) and mean absolute error (MAE) to evaluate forecasting accuracy. Given the ground-truth values $y_i$ and their predictions $\widehat y_i$, the two metrics are defined as

\[
\operatorname{MSE}
=
\frac{1}{N}
\sum_{i=1}^{N}
\left(y_i-\widehat y_i\right)^2,
\qquad
\operatorname{MAE}
=
\frac{1}{N}
\sum_{i=1}^{N}
\left|y_i-\widehat y_i\right|,
\]


where $N$ is the total number of scalar predictions over all evaluation windows, forecast steps and variables for the requested horizon. Lower MSE and MAE indicate more accurate forecasts.

\subsection{Datasets}


We evaluate ISCF-BSCA on seven widely used multivariate time-series forecasting datasets spanning electricity transformers, electricity consumption, meteorology and solar generation. These benchmarks differ substantially in variable dimension, sampling frequency and temporal dynamics. Table~\ref{tab:dataset-statistics} summarizes their basic statistics.

\begin{enumerate}
\item \textbf{ETT (Electricity Transformer Temperature)} records seven load and temperature factors from two electricity transformers in two regions of China between 2016 and 2018 \citep{zhou2021informer}. We use its four standard subsets: ETTh1 and ETTh2 are sampled hourly, whereas ETTm1 and ETTm2 are sampled every 15 minutes.
\item \textbf{Electricity (ECL)} records the hourly electricity consumption of 321 clients from 2012 to 2014 and is provided by the UCI Machine Learning Repository \citep{wu2023timesnet}.
\item \textbf{Weather} is collected by the Beutenberg Weather Station at the Max Planck Institute for Biogeochemistry in Jena, Germany. It contains 21 meteorological indicators sampled every 10 minutes during 2020 \citep{wu2023timesnet}.
\item \textbf{Solar} contains solar-power measurements collected every 10 minutes from 137 photovoltaic plants in Alabama during 2006 \citep{liu2024itransformer}.
\end{enumerate}

\begin{table*}[t]
\centering
\scriptsize
\setlength{\tabcolsep}{4pt}
\caption{Dataset Statistics.}
\label{tab:dataset-statistics}
\resizebox{\textwidth}{!}{%
\begin{tabular}{lcccc}
\toprule
Dataset & Variables & Sampling Frequency & Dataset Size & Domain \\
\midrule
ETTh1 & 7 & 1 h & (8545, 2881, 2881) & Electricity \\
ETTh2 & 7 & 1 h & (8545, 2881, 2881) & Electricity \\
ETTm1 & 7 & 15 min & (34465, 11521, 11521) & Electricity \\
ETTm2 & 7 & 15 min & (34465, 11521, 11521) & Electricity \\
ECL & 321 & 1 h & (18317, 2633, 5261) & Electricity \\
Weather & 21 & 10 min & (36792, 5271, 10540) & Weather \\
Solar & 137 & 10 min & (36601, 5161, 10417) & Solar energy \\
\bottomrule
\end{tabular}%
}
\end{table*}


Dataset Size is reported as (Train, Validation, Test). The ETT datasets follow their standard fixed chronological splits, whereas Weather, ECL and Solar use chronological 70/10/20 splits. The scaler is fitted exclusively on the training segment. All datasets are evaluated at prediction lengths $H\in\{96,192,336,720\}$.

\subsection{Implementation Details}


The local experiments are implemented in Python 3.12.13 and PyTorch 2.9.0 with CUDA 12.8. Each training run uses one NVIDIA GeForce RTX 3090 GPU. ISCF-BSCA is trained from scratch with AdamW and a cosine learning-rate schedule. Table~\ref{tab:training-settings} reports the dataset-specific optimization settings, and Table~\ref{tab:model-settings} lists the corresponding Encoder interface and ISCF-BSCA configuration.

\begin{table*}[t]
\centering
\scriptsize
\setlength{\tabcolsep}{4pt}
\caption{Training and Evaluation Settings.}
\label{tab:training-settings}
\resizebox{\textwidth}{!}{%
\begin{tabular}{lcccccc}
\toprule
Dataset & Max Length & Initial LR & Batch & Grad. Accum. & Max epochs & Patience \\
\midrule
ETTh1 & 720 & $3\times10^{-4}$ & 32 & 1 & 45 & 10 \\
ETTh2 & 720 & $5\times10^{-4}$ & 32 & 1 & 30 & 7 \\
ETTm1 & 720 & $1\times10^{-4}$ & 16 & 2 & 30 & 7 \\
ETTm2 & 720 & $5\times10^{-5}$ & 16 & 2 & 60 & 12 \\
Weather & 720 & $2\times10^{-5}$ & 32 & 1 & 120 & 24 \\
ECL & 720 & $5\times10^{-4}$ & 4 & 8 & 30 & 7 \\
Solar & 720 & $3\times10^{-4}$ & 16 & 2 & 45 & 10 \\
\bottomrule
\end{tabular}%
}
\end{table*}


Gradient Accumulation Steps denotes the number of mini-batches whose gradients are accumulated before one optimizer update; the effective batch size is therefore the batch size multiplied by this value. All runs use seed 2021. The checkpoint with the lowest validation mean MSE over $H\in\{96,192,336,720\}$ is retained and serves all four forecasting requests. Weight decay is $0.01$ except for ETTm2, where it is $0.001$, and early stopping uses zero minimum improvement.

\begin{table*}[t]
\centering
\scriptsize
\setlength{\tabcolsep}{4pt}
\caption{ISCF-BSCA Configuration.}
\label{tab:model-settings}
\resizebox{\textwidth}{!}{%
\begin{tabular}{lccccccccc}
\toprule
Dataset & Scope set & Mode rank & Patch number & $d_{\mathrm{model}}$ & $d_{\mathrm{ff}}$ & Dropout & Weight decay & $\lambda_{\mathrm{scope}}$ & $\lambda_{\mathrm{balance}}^{\max}$ / ramp \\
\midrule
ETTh1 & $\{1,48,144,360,720\}$ & 109 & 24 & 32 & 32 & 0.1 & 0.01 & 1.0 & 0.1 / 0.25 \\
ETTh2 & $\{1,48,144,360,720\}$ & 116 & 12 & 64 & 128 & 0.1 & 0.01 & 1.0 & 0.1 / 0.25 \\
ETTm1 & $\{1,48,144,360,720\}$ & 116 & 1 & 128 & 256 & 0.9 & 0.01 & 1.0 & 0.1 / 0.25 \\
ETTm2 & $\{1,48,144,360,720\}$ & 64 & 6 & 128 & 128 & 0.2 & 0.001 & 1.0 & 0.1 / 0.25 \\
Weather & $\{1,48,144,360,720\}$ & 116 & 19 & 64 & 128 & 0.0 & 0.01 & 1.0 & 0.1 / 0.25 \\
ECL & $\{1,48,144,360,720\}$ & 64 & 1 & 256 & 1,024 & 0.3 & 0.01 & 1.0 & 0.1 / 0.25 \\
Solar & $\{1,48,144,360,720\}$ & 128 & 4 & 256 & 256 & 0.3 & 0.01 & 1.0 & 0.1 / 0.25 \\
\bottomrule
\end{tabular}%
}
\end{table*}


The scope set and BSCA coefficients are shared across datasets. Specifically, $\lambda_{\mathrm{scope}}=1$ and $\lambda_{\mathrm{balance}}(u)=0.1\min(u/0.25,1)$, where $u\in[0,1]$ denotes normalized optimizer progress.

\setcounter{table}{0}
\setcounter{figure}{0}
\section{Full Results}
\label{app:full-results}


Tables~\ref{tab:appendix-horizon-specific} and~\ref{tab:appendix-one-model} provide the complete results for the two forecasting comparisons in Sections~\ref{sec:horizon-specific-comparison} and~\ref{sec:one-model-comparison}. The main text reports the four-horizon average for each dataset, whereas the appendix retains MSE and MAE at every evaluated horizon. Table~\ref{tab:appendix-horizon-specific} compares one unified ISCF-BSCA model with baselines optimized separately for each requested horizon. Table~\ref{tab:appendix-one-model} places every method under the same one-model-all-horizons workflow, where shorter requests are evaluated from the corresponding prefixes of one maximum-length forecast.

\clearpage
% Required packages: booktabs, multirow, graphicx, xcolor
\begin{landscape}
\begin{table}[p]
\centering
\scriptsize
\setlength{\tabcolsep}{1.2pt}
\resizebox{0.76\linewidth}{!}{%
\begin{tabular}{cc|cc|cc|cc|cc|cc|cc|cc|cc|cc|cc|cc}
\toprule
Dataset & H & \multicolumn{2}{c}{\shortstack{HoriScope\\(Ours)}} & \multicolumn{2}{c}{\shortstack{TimeAlign\\(2026)}} & \multicolumn{2}{c}{\shortstack{QDF\\(2026)}} & \multicolumn{2}{c}{\shortstack{AMD\\(2025)}} & \multicolumn{2}{c}{\shortstack{SimpleTM\\(2025)}} & \multicolumn{2}{c}{\shortstack{TVNet\\(2025)}} & \multicolumn{2}{c}{\shortstack{iTransformer\\(2024)}} & \multicolumn{2}{c}{\shortstack{TimeMixer\\(2024)}} & \multicolumn{2}{c}{\shortstack{ModernTCN\\(2024)}} & \multicolumn{2}{c}{\shortstack{PatchTST\\(2023)}} & \multicolumn{2}{c}{\shortstack{DLinear\\(2023)}} \\
 &  & MSE & MAE & MSE & MAE & MSE & MAE & MSE & MAE & MSE & MAE & MSE & MAE & MSE & MAE & MSE & MAE & MSE & MAE & MSE & MAE & MSE & MAE \\
\midrule
\multirow{5}{*}{ETTm1} & 96 & \textcolor{red}{\textbf{0.265}} & \textcolor{red}{\textbf{0.326}} & \textcolor{blue}{\underline{0.280}} & \textcolor{blue}{\underline{0.331}} & 0.283 & 0.336 & 0.289 & 0.343 & 0.325 & 0.363 & 0.288 & 0.343 & 0.300 & 0.353 & 0.293 & 0.345 & 0.292 & 0.346 & 0.293 & 0.346 & 0.300 & 0.345 \\
 & 192 & \textcolor{red}{\textbf{0.304}} & \textcolor{red}{\textbf{0.349}} & \textcolor{blue}{\underline{0.323}} & \textcolor{blue}{\underline{0.357}} & 0.336 & 0.369 & 0.329 & 0.366 & 0.361 & 0.381 & 0.326 & 0.367 & 0.345 & 0.382 & 0.335 & 0.372 & 0.332 & 0.368 & 0.333 & 0.370 & 0.336 & 0.366 \\
 & 336 & \textcolor{red}{\textbf{0.343}} & \textcolor{blue}{\underline{0.371}} & \textcolor{blue}{\underline{0.352}} & 0.376 & 0.369 & 0.390 & 0.365 & 0.386 & 0.391 & 0.404 & 0.365 & 0.391 & 0.374 & 0.398 & 0.368 & 0.386 & 0.365 & 0.391 & 0.369 & \textcolor{red}{\textbf{0.369}} & 0.367 & 0.386 \\
 & 720 & \textcolor{red}{\textbf{0.404}} & \textcolor{red}{\textbf{0.405}} & \textcolor{blue}{\underline{0.408}} & \textcolor{blue}{\underline{0.409}} & 0.423 & 0.424 & 0.423 & 0.417 & 0.454 & 0.438 & 0.412 & 0.413 & 0.429 & 0.430 & 0.426 & 0.417 & 0.416 & 0.417 & 0.416 & 0.420 & 0.419 & 0.416 \\
\cmidrule(lr){2-24}
 & Avg. & \textcolor{red}{\textbf{0.329}} & \textcolor{red}{\textbf{0.363}} & \textcolor{blue}{\underline{0.341}} & \textcolor{blue}{\underline{0.368}} & 0.353 & 0.380 & 0.351 & 0.378 & 0.383 & 0.396 & 0.348 & 0.379 & 0.362 & 0.391 & 0.356 & 0.380 & 0.351 & 0.381 & 0.353 & 0.376 & 0.356 & 0.378 \\
\midrule
\multirow{5}{*}{ETTm2} & 96 & \textcolor{blue}{\underline{0.159}} & \textcolor{red}{\textbf{0.242}} & \textcolor{red}{\textbf{0.157}} & \textcolor{blue}{\underline{0.244}} & 0.164 & 0.253 & 0.168 & 0.258 & 0.175 & 0.258 & 0.161 & 0.254 & 0.175 & 0.266 & 0.165 & 0.256 & 0.166 & 0.256 & 0.166 & 0.256 & 0.164 & 0.255 \\
 & 192 & \textcolor{blue}{\underline{0.218}} & \textcolor{blue}{\underline{0.285}} & \textcolor{red}{\textbf{0.210}} & \textcolor{red}{\textbf{0.280}} & \textcolor{blue}{\underline{0.218}} & 0.291 & 0.221 & 0.295 & 0.240 & 0.301 & 0.220 & 0.293 & 0.242 & 0.312 & 0.225 & 0.298 & 0.222 & 0.293 & 0.223 & 0.296 & 0.224 & 0.304 \\
 & 336 & \textcolor{blue}{\underline{0.270}} & \textcolor{blue}{\underline{0.318}} & \textcolor{red}{\textbf{0.265}} & \textcolor{blue}{\underline{0.318}} & 0.282 & 0.330 & \textcolor{blue}{\underline{0.270}} & 0.327 & 0.305 & 0.342 & 0.272 & \textcolor{red}{\textbf{0.316}} & 0.282 & 0.340 & 0.277 & 0.332 & 0.272 & 0.324 & 0.274 & 0.329 & 0.277 & 0.337 \\
 & 720 & \textcolor{red}{\textbf{0.343}} & \textcolor{red}{\textbf{0.371}} & \textcolor{blue}{\underline{0.346}} & \textcolor{blue}{\underline{0.374}} & 0.364 & 0.385 & 0.355 & 0.381 & 0.401 & 0.399 & 0.349 & 0.379 & 0.378 & 0.398 & 0.360 & 0.387 & 0.351 & 0.381 & 0.362 & 0.385 & 0.371 & 0.401 \\
\cmidrule(lr){2-24}
 & Avg. & \textcolor{blue}{\underline{0.248}} & \textcolor{red}{\textbf{0.304}} & \textcolor{red}{\textbf{0.245}} & \textcolor{red}{\textbf{0.304}} & 0.257 & 0.315 & 0.254 & 0.315 & 0.280 & 0.325 & 0.251 & \textcolor{blue}{\underline{0.311}} & 0.269 & 0.329 & 0.257 & 0.318 & 0.253 & 0.314 & 0.256 & 0.317 & 0.259 & 0.324 \\
\midrule
\multirow{5}{*}{ETTh1} & 96 & \textcolor{red}{\textbf{0.348}} & \textcolor{red}{\textbf{0.386}} & 0.380 & 0.407 & 0.370 & 0.394 & 0.371 & 0.399 & \textcolor{blue}{\underline{0.367}} & \textcolor{blue}{\underline{0.393}} & 0.371 & 0.408 & 0.386 & 0.405 & 0.372 & 0.401 & 0.368 & 0.394 & 0.377 & 0.397 & 0.379 & 0.403 \\
 & 192 & \textcolor{red}{\textbf{0.377}} & \textcolor{red}{\textbf{0.403}} & 0.406 & 0.419 & 0.422 & 0.427 & 0.403 & 0.420 & 0.429 & 0.427 & \textcolor{blue}{\underline{0.398}} & \textcolor{blue}{\underline{0.409}} & 0.424 & 0.440 & 0.413 & 0.430 & 0.405 & 0.413 & 0.409 & 0.428 & 0.408 & 0.419 \\
 & 336 & \textcolor{blue}{\underline{0.393}} & 0.414 & 0.427 & 0.432 & 0.464 & 0.455 & 0.423 & 0.432 & 0.447 & 0.445 & 0.401 & \textcolor{red}{\textbf{0.409}} & 0.449 & 0.460 & 0.438 & 0.450 & \textcolor{red}{\textbf{0.391}} & \textcolor{blue}{\underline{0.412}} & 0.431 & 0.444 & 0.440 & 0.440 \\
 & 720 & \textcolor{red}{\textbf{0.443}} & \textcolor{red}{\textbf{0.458}} & 0.459 & 0.464 & 0.499 & 0.493 & 0.452 & 0.461 & 0.470 & 0.465 & 0.458 & \textcolor{blue}{\underline{0.459}} & 0.495 & 0.487 & 0.486 & 0.484 & \textcolor{blue}{\underline{0.450}} & 0.461 & 0.457 & 0.477 & 0.471 & 0.493 \\
\cmidrule(lr){2-24}
 & Avg. & \textcolor{red}{\textbf{0.390}} & \textcolor{red}{\textbf{0.415}} & 0.418 & 0.431 & 0.439 & 0.442 & 0.412 & 0.428 & 0.428 & 0.433 & 0.407 & 0.421 & 0.439 & 0.448 & 0.427 & 0.441 & \textcolor{blue}{\underline{0.404}} & \textcolor{blue}{\underline{0.420}} & 0.419 & 0.437 & 0.425 & 0.439 \\
\midrule
\multirow{5}{*}{ETTh2} & 96 & \textcolor{red}{\textbf{0.241}} & \textcolor{red}{\textbf{0.314}} & 0.270 & 0.331 & 0.283 & 0.342 & 0.279 & 0.343 & 0.285 & 0.342 & \textcolor{blue}{\underline{0.263}} & \textcolor{blue}{\underline{0.329}} & 0.297 & 0.348 & 0.281 & 0.351 & \textcolor{blue}{\underline{0.263}} & 0.332 & 0.274 & 0.337 & 0.289 & 0.353 \\
 & 192 & \textcolor{red}{\textbf{0.282}} & \textcolor{red}{\textbf{0.345}} & 0.334 & 0.374 & 0.359 & 0.391 & 0.363 & 0.397 & 0.357 & 0.387 & \textcolor{blue}{\underline{0.319}} & \textcolor{blue}{\underline{0.372}} & 0.371 & 0.403 & 0.349 & 0.387 & 0.320 & 0.374 & 0.348 & 0.384 & 0.383 & 0.418 \\
 & 336 & \textcolor{red}{\textbf{0.310}} & \textcolor{red}{\textbf{0.369}} & 0.376 & 0.405 & 0.361 & 0.401 & 0.381 & 0.419 & 0.369 & 0.402 & \textcolor{blue}{\underline{0.311}} & \textcolor{blue}{\underline{0.373}} & 0.404 & 0.428 & 0.366 & 0.413 & 0.313 & 0.376 & 0.377 & 0.416 & 0.448 & 0.465 \\
 & 720 & \textcolor{red}{\textbf{0.392}} & \textcolor{red}{\textbf{0.429}} & 0.406 & 0.436 & 0.407 & 0.440 & 0.442 & 0.467 & 0.414 & 0.436 & \textcolor{blue}{\underline{0.401}} & 0.434 & 0.424 & 0.444 & \textcolor{blue}{\underline{0.401}} & 0.436 & \textcolor{red}{\textbf{0.392}} & \textcolor{blue}{\underline{0.433}} & 0.406 & 0.441 & 0.605 & 0.551 \\
\cmidrule(lr){2-24}
 & Avg. & \textcolor{red}{\textbf{0.306}} & \textcolor{red}{\textbf{0.364}} & 0.347 & 0.387 & 0.353 & 0.394 & 0.366 & 0.406 & 0.356 & 0.392 & 0.324 & \textcolor{blue}{\underline{0.377}} & 0.374 & 0.406 & 0.349 & 0.397 & \textcolor{blue}{\underline{0.322}} & 0.379 & 0.351 & 0.395 & 0.431 & 0.447 \\
\midrule
\multirow{5}{*}{Weather} & 96 & \textcolor{red}{\textbf{0.140}} & \textcolor{red}{\textbf{0.180}} & \textcolor{blue}{\underline{0.143}} & \textcolor{blue}{\underline{0.182}} & 0.152 & 0.202 & 0.148 & 0.203 & 0.157 & 0.204 & 0.147 & 0.198 & 0.159 & 0.208 & 0.147 & 0.198 & 0.149 & 0.200 & 0.149 & 0.198 & 0.170 & 0.230 \\
 & 192 & \textcolor{red}{\textbf{0.181}} & \textcolor{red}{\textbf{0.223}} & \textcolor{blue}{\underline{0.185}} & \textcolor{blue}{\underline{0.224}} & 0.205 & 0.251 & 0.193 & 0.243 & 0.205 & 0.248 & 0.194 & 0.238 & 0.200 & 0.248 & 0.192 & 0.243 & 0.196 & 0.245 & 0.194 & 0.241 & 0.216 & 0.273 \\
 & 336 & \textcolor{red}{\textbf{0.231}} & \textcolor{red}{\textbf{0.263}} & \textcolor{blue}{\underline{0.235}} & \textcolor{red}{\textbf{0.263}} & 0.248 & 0.282 & 0.242 & 0.281 & 0.264 & 0.291 & \textcolor{blue}{\underline{0.235}} & \textcolor{blue}{\underline{0.277}} & 0.253 & 0.289 & 0.247 & 0.284 & 0.238 & \textcolor{blue}{\underline{0.277}} & 0.245 & 0.282 & 0.258 & 0.307 \\
 & 720 & \textcolor{red}{\textbf{0.302}} & \textcolor{red}{\textbf{0.313}} & \textcolor{blue}{\underline{0.308}} & \textcolor{blue}{\underline{0.319}} & 0.341 & 0.351 & 0.315 & 0.332 & 0.346 & 0.344 & \textcolor{blue}{\underline{0.308}} & 0.331 & 0.321 & 0.338 & 0.318 & 0.330 & 0.314 & 0.334 & 0.314 & 0.334 & 0.323 & 0.362 \\
\cmidrule(lr){2-24}
 & Avg. & \textcolor{red}{\textbf{0.214}} & \textcolor{red}{\textbf{0.245}} & \textcolor{blue}{\underline{0.218}} & \textcolor{blue}{\underline{0.247}} & 0.237 & 0.271 & 0.225 & 0.265 & 0.243 & 0.271 & 0.221 & 0.261 & 0.233 & 0.271 & 0.226 & 0.264 & 0.224 & 0.264 & 0.226 & 0.264 & 0.242 & 0.293 \\
\midrule
\multirow{5}{*}{ECL} & 96 & \textcolor{red}{\textbf{0.116}} & \textcolor{red}{\textbf{0.213}} & 0.128 & \textcolor{blue}{\underline{0.218}} & \textcolor{blue}{\underline{0.127}} & 0.222 & 0.131 & 0.228 & 0.143 & 0.237 & 0.142 & 0.223 & 0.138 & 0.237 & 0.142 & 0.247 & 0.129 & 0.226 & 0.129 & 0.222 & 0.140 & 0.237 \\
 & 192 & \textcolor{red}{\textbf{0.136}} & \textcolor{red}{\textbf{0.231}} & 0.145 & \textcolor{blue}{\underline{0.235}} & 0.151 & 0.244 & 0.151 & 0.244 & 0.151 & 0.248 & 0.165 & 0.241 & 0.157 & 0.256 & 0.168 & 0.269 & \textcolor{blue}{\underline{0.143}} & 0.239 & 0.147 & 0.240 & 0.152 & 0.249 \\
 & 336 & \textcolor{red}{\textbf{0.157}} & \textcolor{red}{\textbf{0.251}} & 0.162 & \textcolor{red}{\textbf{0.251}} & 0.167 & 0.260 & 0.167 & 0.262 & 0.176 & 0.269 & 0.164 & 0.269 & 0.167 & 0.264 & 0.171 & 0.260 & \textcolor{blue}{\underline{0.161}} & \textcolor{blue}{\underline{0.259}} & 0.163 & \textcolor{blue}{\underline{0.259}} & 0.170 & 0.267 \\
 & 720 & 0.195 & \textcolor{blue}{\underline{0.280}} & \textcolor{red}{\textbf{0.190}} & \textcolor{red}{\textbf{0.278}} & 0.223 & 0.314 & 0.200 & 0.292 & 0.201 & 0.293 & \textcolor{red}{\textbf{0.190}} & 0.284 & 0.194 & 0.286 & 0.209 & 0.313 & \textcolor{blue}{\underline{0.191}} & 0.286 & 0.197 & 0.290 & 0.203 & 0.301 \\
\cmidrule(lr){2-24}
 & Avg. & \textcolor{red}{\textbf{0.151}} & \textcolor{red}{\textbf{0.244}} & \textcolor{blue}{\underline{0.156}} & \textcolor{blue}{\underline{0.246}} & 0.167 & 0.260 & 0.162 & 0.257 & 0.168 & 0.262 & 0.165 & 0.254 & 0.164 & 0.261 & 0.173 & 0.272 & \textcolor{blue}{\underline{0.156}} & 0.253 & 0.159 & 0.253 & 0.166 & 0.264 \\
\midrule
\multirow{5}{*}{Solar} & 96 & \textcolor{red}{\textbf{0.166}} & \textcolor{red}{\textbf{0.195}} & 0.182 & \textcolor{blue}{\underline{0.205}} & 0.193 & 0.232 & 0.188 & 0.235 & \textcolor{blue}{\underline{0.169}} & 0.229 & 0.204 & 0.260 & 0.188 & 0.232 & 0.179 & 0.232 & 0.202 & 0.263 & 0.178 & 0.229 & 0.197 & 0.210 \\
 & 192 & \textcolor{red}{\textbf{0.185}} & \textcolor{red}{\textbf{0.206}} & 0.198 & \textcolor{blue}{\underline{0.217}} & 0.207 & 0.261 & 0.203 & 0.246 & \textcolor{blue}{\underline{0.189}} & 0.252 & 0.227 & 0.272 & 0.193 & 0.248 & 0.201 & 0.259 & 0.223 & 0.279 & \textcolor{blue}{\underline{0.189}} & 0.246 & 0.218 & 0.222 \\
 & 336 & \textcolor{blue}{\underline{0.195}} & \textcolor{red}{\textbf{0.213}} & 0.203 & \textcolor{blue}{\underline{0.222}} & 0.210 & 0.263 & 0.213 & 0.252 & 0.199 & 0.262 & 0.241 & 0.288 & 0.203 & 0.249 & \textcolor{red}{\textbf{0.190}} & 0.256 & 0.241 & 0.292 & 0.198 & 0.249 & 0.234 & 0.231 \\
 & 720 & 0.204 & \textcolor{red}{\textbf{0.218}} & \textcolor{red}{\textbf{0.201}} & \textcolor{blue}{\underline{0.223}} & 0.223 & 0.278 & 0.215 & 0.257 & 0.207 & 0.259 & 0.242 & 0.287 & 0.223 & 0.261 & \textcolor{blue}{\underline{0.203}} & 0.261 & 0.247 & 0.292 & 0.209 & 0.256 & 0.243 & 0.241 \\
\cmidrule(lr){2-24}
 & Avg. & \textcolor{red}{\textbf{0.188}} & \textcolor{red}{\textbf{0.208}} & 0.196 & \textcolor{blue}{\underline{0.217}} & 0.208 & 0.258 & 0.205 & 0.248 & \textcolor{blue}{\underline{0.191}} & 0.251 & 0.229 & 0.277 & 0.202 & 0.248 & 0.193 & 0.252 & 0.228 & 0.282 & 0.194 & 0.245 & 0.223 & 0.226 \\
\bottomrule
\end{tabular}%
}
\caption{Full results for the horizon-specific comparison. Results are reported as MSE and MAE for $H\in\{96,192,336,720\}$, and Avg. is the arithmetic mean over the four horizons. HoriScope uses one unified model per dataset, whereas the baselines follow their horizon-specific evaluation protocols. The best and second-best displayed values are highlighted in bold and underlined, respectively.}
\label{tab:appendix-horizon-specific}
\end{table}
\end{landscape}




\clearpage
% Required packages: booktabs, multirow, graphicx, xcolor
\begin{table*}[p]
\centering
\scriptsize
\setlength{\tabcolsep}{1.2pt}
\resizebox{\textwidth}{!}{%
\begin{tabular}{cc|cc|cc|cc|cc|cc|cc|cc|cc}
\toprule
Dataset & H & \multicolumn{2}{c}{HoriScope (Ours)} & \multicolumn{2}{c}{TimeAlign (2026)} & \multicolumn{2}{c}{QDF (2026)} & \multicolumn{2}{c}{AMD (2025)} & \multicolumn{2}{c}{SimpleTM (2025)} & \multicolumn{2}{c}{iTransformer (2024b)} & \multicolumn{2}{c}{PatchTST (2023)} & \multicolumn{2}{c}{DLinear (2023)} \\
 &  & MSE & MAE & MSE & MAE & MSE & MAE & MSE & MAE & MSE & MAE & MSE & MAE & MSE & MAE & MSE & MAE \\
\midrule
\multirow{5}{*}{ETTm1} & 96 & \textcolor{red}{\textbf{0.265}} & \textcolor{red}{\textbf{0.326}} & \textcolor{blue}{\underline{0.290}} & \textcolor{blue}{\underline{0.338}} & 0.292 & 0.348 & 0.294 & 0.348 & 0.326 & 0.368 & 0.346 & 0.382 & 0.298 & 0.350 & 0.301 & 0.346 \\
 & 192 & \textcolor{red}{\textbf{0.304}} & \textcolor{red}{\textbf{0.349}} & \textcolor{blue}{\underline{0.325}} & \textcolor{blue}{\underline{0.358}} & 0.330 & 0.371 & 0.333 & 0.368 & 0.362 & 0.385 & 0.384 & 0.399 & 0.335 & 0.372 & 0.336 & 0.367 \\
 & 336 & \textcolor{red}{\textbf{0.343}} & \textcolor{red}{\textbf{0.371}} & \textcolor{blue}{\underline{0.355}} & \textcolor{blue}{\underline{0.378}} & 0.364 & 0.391 & 0.366 & 0.387 & 0.392 & 0.406 & 0.418 & 0.420 & 0.367 & 0.392 & 0.371 & 0.387 \\
 & 720 & \textcolor{red}{\textbf{0.404}} & \textcolor{red}{\textbf{0.405}} & \textcolor{blue}{\underline{0.408}} & \textcolor{blue}{\underline{0.409}} & 0.423 & 0.424 & 0.423 & 0.417 & 0.454 & 0.438 & 0.487 & 0.456 & 0.415 & 0.421 & 0.427 & 0.423 \\
\cmidrule(lr){2-18}
 & Avg. & \textcolor{red}{\textbf{0.329}} & \textcolor{red}{\textbf{0.363}} & \textcolor{blue}{\underline{0.345}} & \textcolor{blue}{\underline{0.371}} & 0.352 & 0.384 & 0.354 & 0.380 & 0.384 & 0.399 & 0.409 & 0.415 & 0.354 & 0.384 & 0.359 & 0.381 \\
\midrule
\multirow{5}{*}{ETTm2} & 96 & \textcolor{red}{\textbf{0.159}} & \textcolor{red}{\textbf{0.242}} & 0.168 & \textcolor{blue}{\underline{0.253}} & \textcolor{blue}{\underline{0.165}} & 0.255 & 0.170 & 0.260 & 0.186 & 0.270 & 0.190 & 0.277 & 0.169 & 0.258 & 0.189 & 0.289 \\
 & 192 & \textcolor{blue}{\underline{0.218}} & \textcolor{red}{\textbf{0.285}} & \textcolor{red}{\textbf{0.216}} & \textcolor{blue}{\underline{0.286}} & 0.223 & 0.294 & 0.222 & 0.295 & 0.247 & 0.308 & 0.253 & 0.315 & 0.223 & 0.295 & 0.256 & 0.339 \\
 & 336 & \textcolor{blue}{\underline{0.270}} & \textcolor{red}{\textbf{0.318}} & \textcolor{red}{\textbf{0.268}} & \textcolor{blue}{\underline{0.320}} & 0.276 & 0.330 & 0.271 & 0.327 & 0.305 & 0.344 & 0.318 & 0.354 & 0.275 & 0.329 & 0.320 & 0.381 \\
 & 720 & \textcolor{red}{\textbf{0.343}} & \textcolor{red}{\textbf{0.371}} & \textcolor{blue}{\underline{0.346}} & \textcolor{blue}{\underline{0.374}} & 0.364 & 0.385 & 0.355 & 0.381 & 0.401 & 0.399 & 0.416 & 0.409 & 0.365 & 0.383 & 0.403 & 0.428 \\
\cmidrule(lr){2-18}
 & Avg. & \textcolor{red}{\textbf{0.248}} & \textcolor{red}{\textbf{0.304}} & \textcolor{blue}{\underline{0.250}} & \textcolor{blue}{\underline{0.308}} & 0.257 & 0.316 & 0.255 & 0.316 & 0.285 & 0.330 & 0.294 & 0.339 & 0.258 & 0.316 & 0.292 & 0.359 \\
\midrule
\multirow{5}{*}{ETTh1} & 96 & \textcolor{red}{\textbf{0.348}} & \textcolor{red}{\textbf{0.386}} & \textcolor{blue}{\underline{0.375}} & 0.404 & 0.401 & 0.418 & 0.378 & \textcolor{blue}{\underline{0.403}} & 0.382 & 0.406 & 0.409 & 0.422 & 0.394 & 0.414 & 0.416 & 0.436 \\
 & 192 & \textcolor{red}{\textbf{0.377}} & \textcolor{red}{\textbf{0.403}} & 0.411 & 0.425 & 0.441 & 0.441 & \textcolor{blue}{\underline{0.410}} & \textcolor{blue}{\underline{0.424}} & 0.434 & 0.431 & 0.457 & 0.448 & 0.429 & 0.433 & 0.457 & 0.465 \\
 & 336 & \textcolor{red}{\textbf{0.393}} & \textcolor{red}{\textbf{0.414}} & 0.433 & 0.437 & 0.467 & 0.459 & \textcolor{blue}{\underline{0.426}} & \textcolor{blue}{\underline{0.433}} & 0.448 & 0.442 & 0.496 & 0.467 & 0.440 & 0.442 & 0.484 & 0.481 \\
 & 720 & \textcolor{red}{\textbf{0.443}} & \textcolor{red}{\textbf{0.458}} & 0.459 & 0.464 & 0.499 & 0.493 & \textcolor{blue}{\underline{0.452}} & \textcolor{blue}{\underline{0.461}} & 0.470 & 0.465 & 0.509 & 0.494 & 0.455 & 0.469 & 0.538 & 0.538 \\
\cmidrule(lr){2-18}
 & Avg. & \textcolor{red}{\textbf{0.390}} & \textcolor{red}{\textbf{0.415}} & 0.420 & 0.433 & 0.452 & 0.453 & \textcolor{blue}{\underline{0.416}} & \textcolor{blue}{\underline{0.430}} & 0.434 & 0.436 & 0.468 & 0.458 & 0.429 & 0.440 & 0.474 & 0.480 \\
\midrule
\multirow{5}{*}{ETTh2} & 96 & \textcolor{red}{\textbf{0.241}} & \textcolor{red}{\textbf{0.314}} & 0.296 & 0.349 & \textcolor{blue}{\underline{0.272}} & 0.340 & 0.285 & 0.351 & 0.285 & 0.344 & 0.303 & 0.353 & 0.277 & \textcolor{blue}{\underline{0.339}} & 0.322 & 0.382 \\
 & 192 & \textcolor{red}{\textbf{0.282}} & \textcolor{red}{\textbf{0.345}} & 0.364 & 0.393 & \textcolor{blue}{\underline{0.334}} & \textcolor{blue}{\underline{0.382}} & 0.374 & 0.405 & 0.358 & 0.388 & 0.385 & 0.401 & 0.345 & \textcolor{blue}{\underline{0.382}} & 0.367 & 0.408 \\
 & 336 & \textcolor{red}{\textbf{0.310}} & \textcolor{red}{\textbf{0.369}} & 0.405 & 0.418 & \textcolor{blue}{\underline{0.358}} & \textcolor{blue}{\underline{0.404}} & 0.379 & 0.421 & 0.367 & \textcolor{blue}{\underline{0.404}} & 0.424 & 0.433 & 0.379 & 0.414 & 0.385 & 0.425 \\
 & 720 & \textcolor{red}{\textbf{0.392}} & \textcolor{red}{\textbf{0.429}} & \textcolor{blue}{\underline{0.406}} & \textcolor{blue}{\underline{0.436}} & 0.407 & 0.440 & 0.442 & 0.467 & 0.414 & \textcolor{blue}{\underline{0.436}} & 0.430 & 0.446 & 0.407 & 0.443 & 0.584 & 0.543 \\
\cmidrule(lr){2-18}
 & Avg. & \textcolor{red}{\textbf{0.306}} & \textcolor{red}{\textbf{0.364}} & 0.368 & 0.399 & \textcolor{blue}{\underline{0.343}} & \textcolor{blue}{\underline{0.392}} & 0.370 & 0.411 & 0.356 & 0.393 & 0.385 & 0.408 & 0.352 & 0.395 & 0.415 & 0.439 \\
\midrule
\multirow{5}{*}{Weather} & 96 & \textcolor{red}{\textbf{0.140}} & \textcolor{red}{\textbf{0.180}} & \textcolor{blue}{\underline{0.148}} & \textcolor{blue}{\underline{0.187}} & 0.191 & 0.246 & 0.149 & 0.204 & 0.177 & 0.222 & 0.179 & 0.222 & 0.155 & 0.205 & 0.177 & 0.240 \\
 & 192 & \textcolor{red}{\textbf{0.181}} & \textcolor{red}{\textbf{0.223}} & \textcolor{blue}{\underline{0.187}} & \textcolor{blue}{\underline{0.226}} & 0.231 & 0.276 & 0.194 & 0.244 & 0.221 & 0.259 & 0.228 & 0.262 & 0.199 & 0.246 & 0.218 & 0.278 \\
 & 336 & \textcolor{red}{\textbf{0.231}} & \textcolor{red}{\textbf{0.263}} & \textcolor{blue}{\underline{0.238}} & \textcolor{blue}{\underline{0.268}} & 0.276 & 0.307 & 0.244 & 0.282 & 0.273 & 0.296 & 0.283 & 0.301 & 0.249 & 0.284 & 0.264 & 0.317 \\
 & 720 & \textcolor{red}{\textbf{0.302}} & \textcolor{red}{\textbf{0.313}} & \textcolor{blue}{\underline{0.308}} & \textcolor{blue}{\underline{0.319}} & 0.341 & 0.351 & 0.315 & 0.332 & 0.346 & 0.344 & 0.359 & 0.351 & 0.319 & 0.335 & 0.328 & 0.369 \\
\cmidrule(lr){2-18}
 & Avg. & \textcolor{red}{\textbf{0.214}} & \textcolor{red}{\textbf{0.245}} & \textcolor{blue}{\underline{0.220}} & \textcolor{blue}{\underline{0.250}} & 0.260 & 0.295 & 0.226 & 0.266 & 0.254 & 0.280 & 0.262 & 0.284 & 0.231 & 0.268 & 0.247 & 0.301 \\
\midrule
\multirow{5}{*}{ECL} & 96 & \textcolor{red}{\textbf{0.116}} & \textcolor{red}{\textbf{0.213}} & \textcolor{blue}{\underline{0.129}} & \textcolor{blue}{\underline{0.220}} & 0.155 & 0.258 & 0.132 & 0.230 & 0.150 & 0.246 & 0.152 & 0.245 & 0.132 & 0.228 & 0.140 & 0.238 \\
 & 192 & \textcolor{red}{\textbf{0.136}} & \textcolor{red}{\textbf{0.231}} & \textcolor{blue}{\underline{0.145}} & \textcolor{blue}{\underline{0.235}} & 0.169 & 0.269 & 0.152 & 0.246 & 0.152 & 0.251 & 0.164 & 0.257 & 0.148 & 0.243 & 0.154 & 0.251 \\
 & 336 & \textcolor{red}{\textbf{0.157}} & \textcolor{red}{\textbf{0.251}} & \textcolor{blue}{\underline{0.163}} & \textcolor{blue}{\underline{0.252}} & 0.185 & 0.284 & 0.166 & 0.261 & 0.173 & 0.269 & 0.179 & 0.273 & 0.164 & 0.259 & 0.169 & 0.268 \\
 & 720 & \textcolor{blue}{\underline{0.195}} & \textcolor{blue}{\underline{0.280}} & \textcolor{red}{\textbf{0.190}} & \textcolor{red}{\textbf{0.278}} & 0.223 & 0.314 & 0.200 & 0.292 & 0.201 & 0.293 & 0.211 & 0.301 & 0.201 & 0.291 & 0.204 & 0.301 \\
\cmidrule(lr){2-18}
 & Avg. & \textcolor{red}{\textbf{0.151}} & \textcolor{red}{\textbf{0.244}} & \textcolor{blue}{\underline{0.157}} & \textcolor{blue}{\underline{0.246}} & 0.183 & 0.281 & 0.163 & 0.257 & 0.169 & 0.265 & 0.177 & 0.269 & 0.161 & 0.255 & 0.167 & 0.265 \\
\midrule
\multirow{5}{*}{Solar} & 96 & \textcolor{red}{\textbf{0.166}} & \textcolor{red}{\textbf{0.195}} & 0.179 & \textcolor{blue}{\underline{0.211}} & 0.186 & 0.256 & 0.190 & 0.242 & \textcolor{blue}{\underline{0.171}} & 0.233 & 0.197 & 0.239 & 0.181 & 0.235 & 0.222 & 0.292 \\
 & 192 & \textcolor{red}{\textbf{0.185}} & \textcolor{red}{\textbf{0.206}} & 0.192 & \textcolor{blue}{\underline{0.218}} & 0.205 & 0.267 & 0.202 & 0.249 & \textcolor{blue}{\underline{0.187}} & 0.246 & 0.228 & 0.259 & 0.190 & 0.245 & 0.250 & 0.309 \\
 & 336 & \textcolor{blue}{\underline{0.195}} & \textcolor{red}{\textbf{0.213}} & 0.198 & \textcolor{blue}{\underline{0.223}} & 0.215 & 0.272 & 0.210 & 0.254 & \textcolor{red}{\textbf{0.193}} & 0.251 & 0.243 & 0.271 & 0.198 & 0.251 & 0.269 & 0.324 \\
 & 720 & \textcolor{blue}{\underline{0.204}} & \textcolor{red}{\textbf{0.218}} & \textcolor{red}{\textbf{0.201}} & \textcolor{blue}{\underline{0.223}} & 0.223 & 0.278 & 0.215 & 0.257 & 0.207 & 0.259 & 0.250 & 0.276 & 0.206 & 0.257 & 0.271 & 0.327 \\
\cmidrule(lr){2-18}
 & Avg. & \textcolor{red}{\textbf{0.188}} & \textcolor{red}{\textbf{0.208}} & 0.193 & \textcolor{blue}{\underline{0.219}} & 0.207 & 0.268 & 0.204 & 0.250 & \textcolor{blue}{\underline{0.190}} & 0.247 & 0.229 & 0.261 & 0.194 & 0.247 & 0.253 & 0.313 \\
\bottomrule
\end{tabular}%
}
\caption{Full results for the one-model-all-horizons comparison. Each method uses one maximum-horizon model per dataset, and shorter horizons are evaluated from the corresponding output prefixes. Results are reported as MSE and MAE, and Avg. denotes the arithmetic mean over $H\in\{96,192,336,720\}$. The best and second-best displayed values are highlighted in bold and underlined, respectively.}
\label{tab:appendix-one-model}
\end{table*}




\setcounter{table}{0}
\setcounter{figure}{0}
\section{Visualization}
\label{app:visualization}


Fig.~\ref{fig:appendix-visualization} presents supplementary UVHF examples for the seven paper-core datasets. For each dataset, two validation samples show the ground-truth future and the fused ISCF-BSCA prediction over the maximum target length $T=720$. The paired rulers and vertical guides mark the four supported prediction lengths, showing how the same predicted trajectory provides the prefixes requested at $H\in\{96,192,336,720\}$.


\begin{figure}[H]
\centering
\includegraphics[width=0.90\textwidth]{figures/figure_c1_varied_horizon_forecasts.pdf}
\caption{Visualization of varied-horizon forecasts across seven datasets. The dark curve denotes the ground truth and the teal curve denotes the prediction from the frozen unified ISCF-BSCA model. Two validation examples are shown for each dataset. The paired rulers and vertical guides identify the four supported prediction lengths.}
\label{fig:appendix-visualization}
\end{figure}
